\documentclass[11pt,a4paper]{article}

\usepackage[utf8]{inputenc}
\usepackage[T1]{fontenc}
\usepackage{lmodern}
\usepackage[spanish,es-noshorthands,es-tabla,es-nodecimaldot]{babel}
\usepackage{microtype}
\usepackage[margin=2.5cm]{geometry}
\usepackage{amsmath,amssymb,amsthm}
\usepackage{booktabs,longtable,tabularx}
\usepackage{enumitem}
\usepackage{xcolor}
\usepackage{listings}
\usepackage{needspace}
\usepackage{tikz}
\usetikzlibrary{arrows.meta,positioning,shapes.geometric,fit,calc}
\usepackage[hidelinks]{hyperref}

% ---------------------------------------------------------------------------
% Notación
\newcommand{\isa}[1]{\texttt{#1}}
\newcommand{\cmd}[1]{\textbf{\texttt{#1}}}
\newcommand{\R}{\mathcal{R}}
\newcommand{\E}{\mathcal{E}}
\newcommand{\Sset}{\mathcal{S}}
\newcommand{\rw}{\rightarrow}
\newcommand{\rws}{\stackrel{*}{\rightarrow}}
\newcommand{\lpo}{>_{\mathit{lpo}}}

\theoremstyle{definition}
\newtheorem{definicion}{Definición}[section]
\newtheorem{ejemplo}[definicion]{Ejemplo}
\theoremstyle{remark}
\newtheorem*{nota}{Nota}

% ---------------------------------------------------------------------------
% Listados Isar / ML
\definecolor{kw}{rgb}{0.0,0.25,0.55}
\definecolor{cm}{rgb}{0.35,0.35,0.35}
\lstdefinelanguage{Isar}{
  morekeywords={theory,imports,begin,end,datatype,fun,definition,where,lemma,
    by,apply,done,oops,complete,orient_rules,schematic_lemma,declare,thm,ML,
    for,and,assumes,shows,induct_auto,inductive_prover,inductive_provers,
    inductive_sledgehammer_prover,spass_auto,rec_complete,setup},
  sensitive=true,
  morecomment=[s]{(*}{*)},
  morestring=[b]",
}
\lstset{
  language=Isar,
  basicstyle=\ttfamily\small,
  keywordstyle=\color{kw}\bfseries,
  commentstyle=\color{cm}\itshape,
  stringstyle=\color{black},
  columns=fullflexible,
  keepspaces=true,
  showstringspaces=false,
  upquote=true,
  frame=single,
  framesep=4pt,
  xleftmargin=6pt,xrightmargin=6pt,
  breaklines=true,
  extendedchars=true,
  literate=
    {⇒}{{$\Rightarrow$}}1 {⟹}{{$\Longrightarrow$}}2 {≡}{{$\equiv$}}1
    {⋀}{{$\bigwedge$}}1 {λ}{{$\lambda$}}1 {∀}{{$\forall$}}1 {∃}{{$\exists$}}1
    {∧}{{$\wedge$}}1 {∨}{{$\vee$}}1 {¬}{{$\neg$}}1 {⟶}{{$\longrightarrow$}}2
    {‹}{{$\langle$}}1 {›}{{$\rangle$}}1 {≠}{{$\neq$}}1 {→}{{$\rightarrow$}}1
    {á}{{\'a}}1 {é}{{\'e}}1 {í}{{\'i}}1 {ó}{{\'o}}1 {ú}{{\'u}}1 {ñ}{{\~n}}1
}

\title{\textbf{IsaLibs}\\[2pt] Comandos Isar, métodos de prueba\\ y la teoría que hay detrás\\[4pt]
\large Guía técnica para la versión portada a Isabelle2026}
\author{Documentación de IsaScheme / IsaLibs}
\date{Noviembre de 2026}

\begin{document}
\maketitle

\begin{abstract}
IsaLibs es una biblioteca de razonamiento automático para Isabelle/HOL escrita
en Isabelle/ML. Implementa demostración inductiva con descubrimiento de lemas,
compleción de Knuth--Bendix sin fallo (\emph{unfailing}), exploración de teorías
guiada por pruebas aleatorias, críticos de prueba (divergencia de Walsh, reglas de
onda transversales) y comprobación de terminación con el demostrador externo AProVE.
Para cada elemento de la interfaz Isar que exporta la biblioteca (comandos, métodos
de prueba, atributos y opciones de configuración), este documento describe su
sintaxis, lo que hace exactamente según el código fuente y la teoría en la que se
basa. Al
final hay un mapa de módulos ML y una lista de advertencias prácticas.
\end{abstract}

\tableofcontents
\newpage

% ===========================================================================
\section{Introducción}

\subsection{Qué es IsaLibs}

IsaLibs (directorio \isa{IsaLibs/} del repositorio IsaScheme) es la evolución de
\emph{IsaScheme}, un sistema de descubrimiento de teoremas e invención de conceptos
basado en \emph{esquemas}~\cite{montano2012}. A partir de ese núcleo se añadieron:
\begin{itemize}[nosep]
  \item un demostrador inductivo recursivo que analiza sus propios fallos para
        conjeturar lemas (método \cmd{inductive\_prover});
  \item compleción de Knuth--Bendix sin fallo para mantener un conjunto de
        reglas de reescritura confluente y terminante (comando \cmd{complete});
  \item exploración de teorías ecuacionales al estilo QuickSpec, basada en la
        enumeración de $\lambda$-términos y en pruebas con valores concretos;
  \item una interfaz con AProVE para decidir (de forma incompleta) la terminación
        de los sistemas de reescritura que se construyen;
  \item utilidades experimentales (programación genética sobre $\lambda$-cálculo
        tipado, lectura de problemas TIP en SMT-LIB, salida a \LaTeX, gnuplot y MySQL).
\end{itemize}

\subsection{Cómo se carga}
\label{sec:carga}

La sesión \isa{IsaLibs} (archivo \isa{IsaLibs/ROOT}) se basa en \isa{HOL} y en la
teoría \isa{THF\_Arith} de \isa{HOL-TPTP}. La teoría \isa{IsaLibs} está declarada
\emph{global}, así que una teoría de usuario sólo necesita:
\begin{lstlisting}
theory MiTeoria
  imports IsaLibs
begin
  ...
end
\end{lstlisting}
y se procesa con
\begin{lstlisting}[language={}]
isabelle build -d ruta/a/IsaLibs -b IsaLibs        (una vez, construye el heap)
isabelle jedit -d ruta/a/IsaLibs -l IsaLibs MiTeoria.thy
\end{lstlisting}
Para que la terminación funcione hace falta AProVE (sección~\ref{sec:aprove}). Como
\isa{IsaLibs.thy} declara \isa{use\_aprove\_srv = true}, el servidor de AProVE debe
estar corriendo:
\begin{lstlisting}[language={}]
cd ~/AProVE && PATH="$HOME/AProVE/bin:$PATH" java -ea -jar aprove.jar -u srv
\end{lstlisting}
Si AProVE no responde, \emph{todas} las comprobaciones de terminación devuelven
``falso''. Entonces ninguna regla se puede orientar y \cmd{inductive\_prover} no
demuestra ni siquiera metas triviales, sin dar ningún mensaje de error.

\paragraph{Comprobar y reiniciar el servidor.} El servidor puede quedar bloqueado,
sobre todo después de interrumpir una compleción que diverge. En ese estado sigue
escuchando en el puerto 5250, pero contesta \texttt{TIMEOUT} a todo, incluso a
sistemas triviales.
Para comprobarlo se le envía un TRS trivial; la primera línea de la respuesta debe
ser \texttt{YES}:
\begin{lstlisting}[language={}]
printf '(VAR x)\n(RULES f(s(x)) -> f(x))\n' > /tmp/prueba.trs
printf '2\n/tmp/prueba.trs\n5\n' | nc localhost 5250 | head -1
\end{lstlisting}
Si no contesta \texttt{YES}, se detiene el servidor y se lanza de nuevo:
\begin{lstlisting}[language={}]
pkill -f 'aprove.jar -u srv'
cd ~/AProVE && PATH="$HOME/AProVE/bin:$PATH" java -ea -jar aprove.jar -u srv
\end{lstlisting}
\needspace{5\baselineskip}
Para que no ocupe una terminal, la segunda línea puede lanzarse en segundo plano:
\begin{lstlisting}[language={}]
cd ~/AProVE && PATH="$HOME/AProVE/bin:$PATH" \
  nohup java -ea -jar aprove.jar -u srv > server.out 2>&1 &
\end{lstlisting}
Después hay que reiniciar también Isabelle (jEdit), porque
\isa{Aprove.memoized\_terminates} guarda en una tabla global los ``falso'' que dio
el servidor bloqueado, y esas respuestas se seguirían usando mientras viva el
proceso.

\subsection{Resumen de la interfaz Isar}

\begin{longtable}{@{}p{5.6cm}>{\raggedright\arraybackslash}p{2.2cm}>{\raggedright\arraybackslash}p{6.6cm}@{}}
\toprule
\textbf{Nombre} & \textbf{Tipo} & \textbf{Propósito}\\
\midrule
\endhead
\cmd{orient\_rules} & comando & Fija la orientación preferida de las ecuaciones que
  encajan con un patrón (sección~\ref{sec:orient}).\\
\cmd{schematic\_lemma} & comando & Registra un patrón de lema esquemático
  (experimental, sin uso actual; sección~\ref{sec:schematic}).\\
\cmd{complete} & comando (meta) & Demuestra una ecuación y la integra, por compleción,
  en el conjunto \isa{compl} (sección~\ref{sec:complete}).\\
\isa{rec\_complete} & comando & Desactivado; versión antigua de \cmd{complete}.\\
\cmd{divergence\_critic\_method} & comando & Fija el método con el que el crítico de
  divergencia de \cmd{complete} demuestra sus conjeturas (sección~\ref{sec:complete}).\\
\midrule
\cmd{induct\_auto} & método & Búsqueda A* de inducción + \isa{auto}
  (sección~\ref{sec:induct_auto}).\\
\cmd{inductive\_prover} & método & Demostrador inductivo con análisis de fallos y
  descubrimiento de lemas (sección~\ref{sec:inductive_prover}).\\
\cmd{inductive\_sledgehammer\_prover} & método & Igual, con Sledgehammer en el
  demostrador atómico.\\
\cmd{inductive\_provers} & método & Igual, probando primero sin Sledgehammer y luego con él.\\
\cmd{spass\_auto} & método & Sledgehammer + reconstrucción con \isa{metis}.\\
\isa{induct\_auto\_sledgehammer\_\-failure\_tac} & método & El demostrador atómico con
  Sledgehammer, como método independiente.\\
\cmd{exponent\_induct} & método & Inducción sobre la variable del exponente de
  \isa{g \textasciicircum\textasciicircum\ n}, con las demás variables arbitrarias.\\
\midrule
\isa{compl} & atributo / conjunto & Reglas producidas por \cmd{complete}.\\
\isa{critic\_simps} & atributo / conjunto & Lemas extra para la táctica predefinida del
  crítico de divergencia.\\
\isa{critic\_rules} & hecho dinámico & Las reglas de la compleción en curso, visibles
  para el método del crítico.\\
\isa{prop\_scheme} & atributo / conjunto & Esquemas de propiedades (asociatividad,
  conmutatividad, \ldots) para conjeturar.\\
\midrule
\isa{@\{typ\_pat\}}, \isa{@\{term\_pat\}} & antiquotación ML & Tipos y términos con
  variables esquemáticas dentro de código ML.\\
\bottomrule
\end{longtable}

% ===========================================================================
\section{Fundamentos teóricos}
\label{sec:teoria}

Esta sección reúne los conceptos que usan varios comandos a la vez. Las secciones
siguientes remiten aquí.

\subsection{Reescritura de términos}

Un \emph{sistema de reescritura} (TRS) $\R$ es un conjunto de reglas $l \rw r$ en el
que $l$ no es una variable y toda variable de $r$ aparece en $l$. La relación $\rw_\R$
reemplaza una instancia $l\sigma$ de un lado izquierdo por $r\sigma$ en cualquier
posición. Las propiedades que importan son tres.
\begin{itemize}[nosep]
  \item $\R$ es \textbf{terminante} si no existen cadenas infinitas
        $t_1 \rw_\R t_2 \rw_\R \cdots$.
  \item $\R$ es \textbf{confluente} si, siempre que $s \rws t_1$ y $s \rws t_2$,
        existe $u$ con $t_1 \rws u$ y $t_2 \rws u$.
  \item $\R$ es \textbf{convergente} si es terminante y confluente. En un TRS
        convergente cada término tiene una forma normal única. Por tanto la
        igualdad ecuacional se decide normalizando ambos lados y comparándolos.
\end{itemize}
Por el lema de Newman, un TRS terminante es confluente si y sólo si es localmente
confluente. A su vez, por el lema de los pares críticos de Knuth--Bendix, es
localmente confluente si y sólo si todos sus \emph{pares críticos} son
\emph{unibles} (\emph{joinable}). Un par crítico surge cuando el lado izquierdo de
una regla unifica con un subtérmino no variable del lado izquierdo de otra:
\[
  l_1 \rw r_1,\quad l_2 \rw r_2,\quad \sigma = \mathrm{mgu}(l_1, l_2|_p)
  \quad\Longrightarrow\quad \langle\, r_2\sigma,\; (l_2\sigma)[r_1\sigma]_p \,\rangle .
\]
La referencia estándar es \cite{baader-nipkow}.

En Isabelle el simplificador hace reescritura condicional de orden superior. Las
reglas que IsaLibs construye se manejan como teoremas $l \equiv r$ (meta-igualdades)
o $l = r$ (igualdades de HOL). Las funciones \isa{Utils.obj\_to\_meta} y
\isa{Utils.meta\_to\_obj} convierten entre ambas formas.

\subsection{Órdenes de reducción y reescritura ordenada}

Para orientar ecuaciones se usan órdenes de simplificación. IsaLibs usa sobre todo
el \textbf{orden lexicográfico de caminos} (LPO) que ofrece Isabelle
(\isa{Term\_Ord.term\_lpo}), parametrizado por una \emph{precedencia} entre
símbolos. Con $f \succ g$ en la precedencia, $s = f(s_1,\ldots,s_m) \lpo t$ si se
cumple alguna de estas condiciones:
\begin{enumerate}[nosep]
  \item $s_i \geq_{\mathit{lpo}} t$ para algún $i$;
  \item $t = g(t_1,\ldots,t_n)$, $f \succ g$ y $s \lpo t_j$ para todo $j$;
  \item $t = f(t_1,\ldots,t_m)$, $s \lpo t_j$ para todo $j$ y
        $(s_1,\ldots,s_m) \lpo^{\mathit{lex}} (t_1,\ldots,t_m)$.
\end{enumerate}

Hay ecuaciones que ninguna orientación vuelve terminantes, como la conmutatividad
$x + y = y + x$. A estas ecuaciones \emph{permutativas} (ambos lados son iguales
salvo un renombramiento de variables, \isa{Utils.permutative\_eq}) se les aplica
\textbf{reescritura ordenada}~\cite{martin-nipkow}, en la que una instancia
$l\sigma \rw r\sigma$ sólo se aplica si $l\sigma \succ r\sigma$ en un orden total sobre términos cerrados.
El simplificador de Isabelle ya hace esto con las reglas permutativas usando su orden
de términos, que IsaLibs puede sustituir con \isa{Simplifier.set\_term\_ord}.

\subsection{Compleción de Knuth--Bendix sin fallo}
\label{sec:kb}

La compleción de Knuth--Bendix~\cite{knuth-bendix} transforma un conjunto de
ecuaciones $\E$ en un TRS convergente $\R$ equivalente. Para ello orienta ecuaciones,
calcula pares críticos y añade los que no son unibles. La versión clásica
\emph{falla} cuando aparece una ecuación que no se puede orientar. La compleción
\emph{sin fallo} (\emph{unfailing completion}) de Bachmair, Dershowitz y
Plaisted~\cite{bdp} conserva esas ecuaciones como ecuaciones y las usa con
reescritura ordenada. Los pares críticos se calculan entonces también con las
ecuaciones, en ambas direcciones (\emph{pares semicríticos}).

Las reglas de inferencia, en la presentación de~\cite{baader-nipkow}, son:
\begin{description}[nosep,font=\normalfont\itshape]
  \item[Deducir] añadir a $\E$ un par (semi)crítico.
  \item[Orientar] mover $s = t$ de $\E$ a $\R$ como $s \rw t$ si $s \succ t$.
  \item[Simplificar] reescribir un lado de una ecuación de $\E$.
  \item[Eliminar] borrar $s = t$ de $\E$ si es trivial o unible.
  \item[Componer] reescribir el lado derecho de una regla.
  \item[Colapsar] si el lado izquierdo de una regla es reducible por otra, devolverla a $\E$.
\end{description}

\paragraph{Unibilidad en términos cerrados (\emph{ground joinability}).}
Con ecuaciones permutativas no basta con comprobar si $s$ y $t$ tienen la misma forma
normal, porque la reescritura ordenada depende de cómo se instancien las variables.
El criterio de Martin y Nipkow~\cite{martin-nipkow} trata las variables de $s = t$
como constantes nuevas. Después considera \emph{todas} las maneras de insertarlas
en la precedencia y, con el LPO inducido por cada una, reescribe de forma ordenada
ambos lados. Si coinciden para toda extensión, $s = t$ es unible para cualquier
instancia cerrada y se puede eliminar. IsaLibs implementa exactamente esta
comprobación en \isa{ground\_joinable}, enumerando las extensiones con
\isa{Orders.extend\_ord\_eval}.

\subsection{Patrones de orden superior}

Los términos de Isabelle son $\lambda$-términos tipados. En general la unificación
de orden superior es indecidible, pero para los \emph{patrones} de Miller~\cite{miller}
es decidible y tiene unificadores más generales únicos. Un patrón es un término en el
que cada variable libre $F$ sólo aparece aplicada a variables ligadas distintas,
$F\,x_1 \ldots x_n$. IsaLibs calcula pares críticos sólo entre reglas cuyos dos lados
izquierdos son patrones (\isa{Pattern.pattern}), y busca subtérminos también bajo
$\lambda$. Los críticos de divergencia (sección~\ref{sec:criticos}) también usan
unificación de patrones (\isa{Pattern.unify}).

\subsection{Terminación automática con AProVE y pares de dependencia}
\label{sec:aprove}

Cada vez que IsaLibs quiere usar un conjunto de reglas para simplificar, primero
comprueba que sea terminante. La función central es
\isa{Aprove.memoized\_terminates}, que memoriza los resultados en una tabla global
cuya clave es el conjunto de reglas ordenado. Funciona así:
\begin{enumerate}[nosep]
  \item Si el lado derecho de alguna regla contiene su lado izquierdo, responde
        \emph{no} de inmediato.
  \item Traduce las reglas al formato WST/TPDB de AProVE
        (\isa{(VAR \ldots) (RULES \ldots)}):
    \begin{itemize}[nosep]
      \item cada constante se escribe por su nombre base (\isa{Long\_Name.base\_name});
      \item los conectivos de HOL pasan a los símbolos \isa{equals}/\isa{iff},
            \isa{and}, \isa{or}, \isa{implies} y \isa{not};
      \item las variables \isa{?x$_i$} se escriben \isa{x$i$};
      \item las fórmulas cuantificadas se rechazan;
      \item las $\lambda$-abstracciones se eliminan por \emph{$\lambda$-lifting},
            que reemplaza cada $\lambda$ por una constante nueva \isa{f} aplicada a
            sus variables libres y añade la regla que la define;
      \item una regla condicional se escribe \isa{l -> r | c1, c2}, y las
            premisas que no son ecuaciones $P$ (o $\neg P$) se convierten antes
            en $P = \mathit{True}$ (o $P = \mathit{False}$).
    \end{itemize}
  \item Llama a AProVE. Con \isa{use\_aprove\_srv = true} envía
        \isa{"2\textbackslash n<archivo>\textbackslash n<tiempo>\textbackslash n"} al
        servidor (\isa{nc \$APROVE\_SERVER \$APROVE\_SERVER\_PORT}). Con
        \isa{false} ejecuta
        \isa{java -jar \$APROVE\_PATH/aprove.jar -m wst -t $T$} a través del guion
        \isa{\$TIMEAOUT\_PATH/timeout}.
        La regla se acepta si la respuesta contiene \isa{YES}.
\end{enumerate}
AProVE~\cite{aprove} combina, entre otras técnicas, el método de \emph{pares de
dependencia}~\cite{arts-giesl}, órdenes de polinomios y de caminos resueltos con SAT,
y criterios de subtérmino.

IsaLibs trae además su propia implementación del método \emph{estático} de pares de
dependencia para sistemas de reescritura de orden superior ``plain function-passing''
(\isa{dependency\_pair.ML}, \cite{kusakari}), con grafo de dependencias, componentes
fuertemente conexas, criterio de subtérmino, filtrado de argumentos y LPO débil
codificado en SAT (\isa{sat\_interface.ML}). En el código actual esta implementación
\emph{no} interviene en \isa{memoized\_terminates}, porque sólo la usa
\isa{Aprove.all\_terminates}, que no se llama desde ningún sitio.

\subsection{Inducción en Isabelle}

IsaLibs usa dos clases de inducción:
\begin{itemize}[nosep]
  \item La inducción \textbf{estructural} se hace sobre una variable libre cuyo
        tipo es un tipo de datos (\isa{Induct.find\_inductT}).
  \item La inducción \textbf{computacional} (de recursión) usa la regla
        \isa{f.induct} que \isa{fun} genera para una función recursiva $f$. Se usa cuando la meta
        contiene un subtérmino $f\,x_1\ldots x_n$ cuyos argumentos son todos
        variables libres. Esta regla sigue exactamente el esquema de recursión de
        $f$, lo que suele dejar hipótesis de inducción más útiles.
\end{itemize}
En ambos casos el resto de las variables libres se puede generalizar
(\isa{arbitrary}) o dejar fijo. Los demostradores prueban las dos opciones.

\subsection{Descubrimiento de lemas}
\label{sec:criticos}

Una prueba inductiva suele atascarse porque hace falta un lema auxiliar.
\cmd{inductive\_prover} combina varias fuentes clásicas de conjeturas:
\begin{description}[nosep,font=\normalfont\bfseries]
  \item[Generalización] Se reemplaza por una variable nueva un subtérmino que
        aparece varias veces, o se eliminan hipótesis, siempre que la conjetura
        resultante no se refute (Boyer y Moore~\cite{boyer-moore}).
  \item[Análisis de fallos] Las submetas en las que se atascó la prueba se
        convierten en conjeturas, como en los \emph{proof critics} de Ireland y
        Bundy~\cite{ireland-bundy}.
  \item[Crítico de divergencia] Si las submetas sucesivas tienen la forma
        $G(U),\ G(H(U)),\ G(H(H(U))),\ldots$, el crítico de Walsh~\cite{walsh}
        propone un lema ``de onda'' de la forma $G(H(U)) = H'(G(U))$, que mueve
        $H$ hacia fuera.
  \item[Reglas de onda transversales] Esta técnica de
        \emph{rippling}~\cite{rippling} sirve para funciones con acumulador. A
        partir de $G(U,\mathit{Acc})$ y $G(H(U),\mathit{Acc})$ se propone
        $G(H(U),\mathit{Acc}) = G(U, F(\mathit{Acc}))$ con $F$ desconocida, y $F$
        se sintetiza demostrando la meta esquemática.
  \item[Exploración de teorías] Como en QuickSpec~\cite{quickspec} e
        IsaCoSy~\cite{isacosy}, se enumeran términos pequeños sobre la signatura
        del problema, se agrupan por su comportamiento con valores de prueba y se
        proponen como ecuaciones los términos de una misma clase.
  \item[Esquemas] Como en IsaScheme~\cite{montano2012}, se instancian plantillas
        de propiedades conocidas (asociatividad, distributividad, \ldots) con las
        funciones del problema.
\end{description}

\subsection{Refutación de conjeturas}

Antes de intentar probar una conjetura, IsaLibs busca un contraejemplo con
Quickcheck~\cite{quickcheck} y, si no lo encuentra, con Nitpick~\cite{nitpick}
(\isa{Counter\_Example.counter\_example}). Los resultados se memorizan en dos redes de
discriminación globales, una de conjeturas refutadas y otra de no refutadas. Así, una
instancia de una conjetura ya refutada se descarta sin volver a llamar a las
herramientas, y lo mismo ocurre con las conjeturas que generalizan a una ya validada.

% ===========================================================================
\section{Comandos Isar}

IsaLibs declara estas palabras clave en \isa{IsaLibs.thy}:
\begin{lstlisting}
keywords "complete" :: thy_goal and
         "orient_rules" :: thy_decl and
         "schematic_lemma" :: thy_decl
\end{lstlisting}

% ---------------------------------------------------------------------------
\subsection{\texttt{orient\_rules}}
\label{sec:orient}

\paragraph{Sintaxis.}
\begin{lstlisting}
orient_rules "patrón_de_ecuación"
\end{lstlisting}
El argumento se lee con \isa{Proof\_Context.read\_term\_pattern}, así que admite
variables esquemáticas \isa{?x} y variables de tipo \isa{?'a}. Debe ser una
ecuación de HOL $l = r$, normalmente con variables de orden superior.

\paragraph{Qué hace.}
Inserta el patrón en una tabla \emph{global} de ML (\isa{Utils.oriented\_rules\_var},
una \isa{Synchronized.var}). La tabla se consulta en dos lugares:
\begin{enumerate}[nosep]
  \item En \isa{Utils.orient\_eq}, que decide la orientación de toda conjetura y de
        todo lema nuevo. Las heurísticas se aplican en este orden:
    \begin{enumerate}[nosep]
      \item si la ecuación encaja (\isa{Pattern.first\_order\_match}) con un patrón de
            \cmd{orient\_rules}, se deja como está; si encaja la ecuación girada,
            se gira;
      \item la \emph{heurística de Lucas} (\isa{definitional\_orientation}) elige
            la orientación $s = t$ si $f_1$ y $f_2$ son los símbolos de cabeza de
            $s$ y $t$, y $f_2$ está entre las dependencias definicionales de $f_1$
            (o $f_2$ no tiene definición);
      \item si la ecuación es una instancia de una ecuación \isa{.simps}
            existente, se orienta igual;
      \item en otro caso, se usa el orden de términos (\isa{ac\_orient}) y se evitan
            las orientaciones ``malas'' (\isa{bad\_equation}), que son las de lado
            derecho trivial, lado izquierdo contenido en el derecho, constructor en
            la cabeza del lado izquierdo o variables nuevas a la derecha.
    \end{enumerate}
  \item En \isa{discriminate\_rules} (compleción básica y condicional). Antes de
        separar reglas terminantes de ecuaciones, cada ecuación se orienta según
        el primer patrón que encaje (\isa{Pattern.matchess}).
\end{enumerate}

\paragraph{Teoría.}
La compleción sólo puede usar una ecuación para reescribir si la orientación
elegida mantiene el sistema terminante. Algunas orientaciones son terminantes por
sí solas pero dejan de serlo al combinarse con otras reglas, y AProVE no puede saber
cuál prefiere el usuario. Con \cmd{orient\_rules} se fija una orientación canónica
para familias enteras de ecuaciones.

\paragraph{Patrones predefinidos.} \isa{IsaLibs.thy} declara estos:
\begin{lstlisting}
orient_rules "?X (?X (?x :: ?'a) (?y :: ?'a)) (?z :: ?'a) = ?X ?x (?X ?y ?z)"
orient_rules "?X (?X (?x :: ?'a) (?y :: ?'a)) (?z :: ?'a) = ?X ?x (?X ?z ?y)"
   (* ... las seis permutaciones de x, y, z ... *)
orient_rules "?X (?x :: ?'a) (?y :: ?'a) = (?x = ?y)"
orient_rules "(?F (?x :: ?'a) = ?F (?y :: ?'a)) = (?x = ?y)"
\end{lstlisting}
Las seis primeras orientan la asociatividad (y sus variantes AC) siempre de
izquierda a derecha, $(x \circ y) \circ z \rw x \circ (y \circ z)$. Por ejemplo,
\isa{Utils.orient\_eq} convierte
\isa{add x (add y z) = add (add x y) z} en
\isa{add (add x y) z = add x (add y z)}. Junto con la
conmutatividad, que se trata como ecuación permutativa, esto da un sistema
AC-terminante. Las dos últimas orientan las caracterizaciones de igualdad y de
inyectividad hacia la forma más simple.

\begin{nota}
La tabla es estado ML global, no un dato de la teoría. Persiste en el \emph{heap} de
la sesión, afecta a todas las teorías cargadas después en el mismo proceso, no se
deshace al salir de un contexto y no hay comando para quitar un patrón.
\end{nota}

% ---------------------------------------------------------------------------
\subsection{\texttt{schematic\_lemma}}
\label{sec:schematic}

\paragraph{Sintaxis.}
\begin{lstlisting}
schematic_lemma "?F (?x :: ?'a) = ?x"
\end{lstlisting}

\paragraph{Qué hace.}
Lee el patrón igual que \cmd{orient\_rules} y lo guarda en otra tabla global,
\isa{Utils.schematic\_lemmas\_var}. La función \isa{Utils.get\_schematic\_lemmas}
devuelve su contenido, pero \textbf{ningún módulo cargado la usa}. El comando quedó
como infraestructura para una exploración de teorías guiada por lemas esquemáticos,
que no llegó a integrarse (el archivo es de junio de 2017). La única funcionalidad
equivalente en activo son las conjeturas por esquemas de \isa{prop\_scheme}
(sección~\ref{sec:prop_scheme}).

% ---------------------------------------------------------------------------
\subsection{\texttt{complete}}
\label{sec:complete}

\paragraph{Sintaxis.}
\begin{lstlisting}
complete "ecuación"
  by método
\end{lstlisting}
Es un comando de meta (\isa{thy\_goal}) que abre una prueba de la proposición dada.
Esa prueba se cierra como cualquier otra (\isa{by induct\_auto},
\isa{by inductive\_prover}, \isa{by simp}, una prueba estructurada, \ldots). No
acepta nombre ni atributos, y el teorema \emph{no} queda con nombre, de modo que
sólo puede llegar al conjunto \isa{compl}.

\paragraph{Qué hace después de la prueba.}
El código (\isa{commands.ML}, \isa{after\_qed}) sigue estos pasos:
\begin{enumerate}
  \item Convierte el teorema demostrado $e$ en meta-igualdad
        (\isa{Utils.obj\_to\_meta}).
  \item Construye el conjunto inicial $E$ con \isa{global\_lemmas}:
    \begin{itemize}[nosep]
      \item calcula la \emph{clausura definicional} de las constantes de $e$, es
            decir, las constantes de $e$, las que aparecen en sus ecuaciones
            \isa{.simps}, y así sucesivamente (\isa{Utils.definitional\_dependencies});
      \item toma todas las ecuaciones \isa{.simps} de esas constantes;
      \item añade las reglas de \isa{compl} que mencionen alguna de ellas.
    \end{itemize}
  \item Ejecuta \isa{Ground\_Completion.run\_completion'} con $E \cup \{e\}$, usando
        \isa{Aprove.memoized\_terminates} como oráculo de terminación.
  \item Calcula y muestra en la traza:
        \begin{itemize}[nosep]
          \item las \textbf{Additions}, que son las reglas del sistema final que
                no estaban en $E$;
          \item las \textbf{Deletions}, que son los elementos de $E$ que ya no
                están en el sistema final porque quedaron redundantes, por ejemplo
                al colapsar.
        \end{itemize}
  \item Añade las \emph{Additions} a \isa{compl} y quita de \isa{compl} las
        \emph{Deletions}. La actualización se hace en la teoría de fondo, en el
        destino local y en el contexto de prueba. Las ecuaciones \isa{.simps} nunca
        se quitan del simplificador, y que aparezcan en ``Deletions'' sólo indica
        que son redundantes en el sistema completado.
\end{enumerate}

\paragraph{El algoritmo de compleción} (\isa{ground\_completion.ML}).
Es la compleción sin fallo de la sección~\ref{sec:kb} y trabaja con tres conjuntos,
$\E$ (ecuaciones por procesar), $\Sset$ (reglas nuevas aún no usadas para generar
pares críticos) y $\R$ (reglas procesadas).
\begin{enumerate}
  \item Al \textbf{inicio}, $\E$ se ordena por tamaño, de modo que primero van los
        términos más pequeños, con más variables y con menor diferencia de
        variables entre lados (\isa{rule\_int\_ord}).
  \item Cada $e \in \E$ se \textbf{orienta} (\isa{orient}) en estos pasos:
    \begin{enumerate}[nosep]
      \item si es posible, se elige la orientación lineal por la izquierda
            (\isa{linear\_if\_possible});
      \item se normaliza con la parte terminante de $\Sset \cup \R$;
      \item se orienta con \isa{Utils.orient\_rule} (las heurísticas de la
            sección~\ref{sec:orient});
      \item si es \emph{unible en términos cerrados} con $\Sset \cup \R$, se descarta;
      \item si es permutativa, se añade tal cual;
      \item si $\{e\} \cup \R$ es terminante según AProVE, se añade orientada; si no,
            se añade girada.
    \end{enumerate}
  \item Al \textbf{añadir una regla} (\isa{add\_rule}) se hace la interreducción. Cada regla de
        $\Sset$ y $\R$ cuyo lado izquierdo reescribe la nueva regla vuelve a $\E$
        (\emph{colapsar}). En las demás se normaliza el lado derecho
        (\emph{componer}).
  \item Para \textbf{elegir} (\isa{choose}) se toma de $\Sset$ la regla mínima según
        \isa{rule\_int\_ord}. Si tiene la forma $c(x) = c(y)$, \isa{heuristic1}
        intenta demostrar la ecuación más fuerte $x = y$ (primero busca un
        contraejemplo; si no lo hay, la prueba con \cmd{induct\_auto}) y, si lo
        logra, la usa en su lugar.
  \item Para \textbf{deducir}, si la opción \isa{generate\_cps} está activa, se calculan los
        pares semicríticos de la regla elegida con $\R$, de $\R$ con la regla y de la
        regla consigo misma. Se usan las cuatro combinaciones de orientación, se
        unifica con patrones de orden superior en posiciones no variables (también
        bajo $\lambda$), y cada par se obtiene como un \emph{teorema} de Isabelle por
        conversiones (\isa{Conv.rewr\_conv}). Los pares pasan a $\E$ y se repite el
        proceso.
  \item Al \textbf{final}, cuando $\Sset$ queda vacío, \isa{make\_irreducible} elimina las
        reglas que son unibles con el resto.
\end{enumerate}
Para la unibilidad en términos cerrados, todas las constantes de la clausura
definicional empiezan con la misma precedencia y luego se prueban todas las
inserciones de las variables (sección~\ref{sec:kb}). Para usar las reglas,
\isa{discriminate\_rules} separa el conjunto en dos partes:
\begin{itemize}[nosep]
  \item las \emph{ecuaciones} permutativas, que se usan con reescritura ordenada;
  \item las \emph{reglas}, que se aceptan de una en una, de menor a mayor tamaño,
        mientras el conjunto siga siendo terminante según AProVE y pase una prueba
        empírica de terminación AC (\isa{ac\_termination}, que reescribe y comprueba
        que no se repite ningún término ni se excede una cota de tamaño).
\end{itemize}
Las reglas incompatibles (``\isa{Incompatible rule}'' en la traza) se quedan fuera
del conjunto de simplificación.

Hay dos detalles más. Primero, \isa{run\_completion} lanza en paralelo dos
compleciones, con el orden de términos en sentido \isa{LESS} y en sentido
\isa{GREATER}, y se queda con la que termina primero (\isa{Par\_List.get\_some}).
Segundo, la opción \isa{generate\_cps} está activa por omisión, así que
\cmd{complete} deduce pares críticos y, si termina, devuelve un sistema
convergente. Con pares críticos la compleción también puede no terminar. Si sólo
interesa orientar, interreducir y eliminar reglas unibles, se declara
\isa{declare [[generate\_cps = false]]}. La misma opción rige la compleción
condicional que usa \cmd{inductive\_prover} (sección~\ref{sec:inductive_prover}).

\paragraph{Corrección.}
Todas las reglas de \isa{compl} son teoremas de Isabelle. La regla inicial viene de
la prueba del usuario, los pares críticos se derivan por conversiones, y las reglas
de \isa{heuristic1} se demuestran con \cmd{induct\_auto}. La unibilidad y la
terminación sólo deciden qué se \emph{descarta} y qué se \emph{usa}, así que un
error en ellas afecta a la completitud o a la eficiencia, pero nunca a la corrección.

\begin{ejemplo}
Con IsaLibs en Isabelle2026 y AProVE:
\begin{lstlisting}
datatype N = Z | S N
fun add :: "N ⇒ N ⇒ N" where
  "add Z y = y" | "add (S x) y = S (add x y)"

complete "add x Z = x"
  by induct_auto
(* Additions: add ?x Z = ?x      Deletions: (ninguna) *)
thm compl
\end{lstlisting}
Y el ejemplo de \isa{IsaLibs\_Examples.thy} (tarda unos 3 minutos):
\begin{lstlisting}
fun app :: "'a list ⇒ 'a list ⇒ 'a list" where ...
fun rev :: "'a list ⇒ 'a list" where
  "rev [] = []" | "rev (x # xs) = app (rev xs) [x]"

declare [[generate_cps = false]]
complete "itrev xs [] = rev xs"
  by inductive_prover
(* Additions: rev ?xs = itrev ?xs [],
              itrev ?xsa [?xa] = app (itrev ?xsa []) [?xa]
   Deletions: rev (?xa # ?xsa) = app (rev ?xsa) [?xa], rev [] = [] *)
\end{lstlisting}
La compleción decidió que la definición eficiente \isa{itrev} sustituye a la
definición original de \isa{rev}.

El ejemplo desactiva los pares críticos porque con ellos esta compleción no
termina. Genera las dos familias infinitas
\begin{align*}
  \mathit{itrev}\ \mathit{xs}\ [x_1,\ldots,x_n] &=
    \mathit{app}\ (\cdots(\mathit{app}\ (\mathit{itrev}\ \mathit{xs}\ [\,])\ [x_1])\cdots)\ [x_n],\\
  \mathit{app}\ (\mathit{itrev}\ \mathit{xs}\ [x_1,\ldots,x_{n-1}])\ [x_n] &=
    \mathit{itrev}\ \mathit{xs}\ [x_1,\ldots,x_n],
\end{align*}
en las que crece una lista de variables distintas y no una torre de un mismo
símbolo unario, así que el crítico de divergencia no las reconoce. La declaración
también desactiva los pares críticos en la compleción condicional de
\cmd{inductive\_prover}.
\end{ejemplo}

{\raggedright
\paragraph{Opciones relevantes.} \isa{generate\_cps}, \isa{divergence\_critic},
\isa{max\_time\_in\_termination}, \isa{use\_aprove}, \isa{use\_aprove\_srv},
\isa{max\_time\_in\_proof} (para \isa{heuristic1}),
\isa{max\_time\_in\_critic\_proof} (para el crítico de divergencia) y
\isa{max\_time\_in\_counter\_ex}.\par}

\paragraph{Crítico de divergencia (prototipo).}
Con pares críticos, la compleción puede no terminar. El caso típico es una familia
infinita de reglas que sólo difieren en la longitud de una torre de una función
unaria. Por ejemplo, la regla $f(g(f\,x)) = f(g\,x)$ genera
$f(g^n(f\,x)) = f(g^n\,x)$ para todo $n \ge 1$. La literatura representa esas
familias con meta-reglas (H.~Kirchner~\cite{kirchner89}), términos de recurrencia
(Chen, Hsiang y Kong~\cite{chen-hsiang-kong}), términos con exponentes
enteros (Comon~\cite{comon95}) o una signatura enriquecida (Thomas y
Watson~\cite{thomas-watson}). En HOL basta la iteración de funciones
\isa{g \textasciicircum\textasciicircum\ n}, porque la familia es \emph{un} teorema,
\isa{f ((g \textasciicircum\textasciicircum\ Suc n) (f x)) = f ((g \textasciicircum\textasciicircum\ Suc n) x)},
que se demuestra por inducción sobre \isa{n}.

\cmd{complete} lleva un crítico (\isa{divergence\_critic}, activo por omisión) que
sólo actúa cuando se deducen pares críticos y hace lo siguiente:
\begin{enumerate}[nosep]
  \item \textbf{Detectar.} Cada regla elegida se abstrae en un \emph{esqueleto},
        donde las torres maximales $c(c(\dots s))$ de símbolos unarios se
        sustituyen por huecos, más la lista de longitudes de esas torres. Si tres
        reglas elegidas tienen el mismo esqueleto y sus longitudes crecen en 0 o 1
        en cada paso (y en alguna crecen), se conjetura la familia.
  \item \textbf{Generalizar.} Las torres que crecen se escriben
        \isa{c \textasciicircum\textasciicircum\ Suc$^k$ n}, donde $k$ es su
        longitud en el miembro más pequeño. Las que no crecen se quedan igual.
  \item \textbf{Demostrar} la regla esquemática. Por omisión, por inducción sobre
        \isa{n} (con las demás variables arbitrarias), simplificación y
        \isa{metis} con las reglas del momento; el método se puede cambiar
        (ver ``Método de prueba del crítico'' más abajo).
  \item \textbf{Comprobar la terminación} con AProVE de las reglas más las
        instancias $n = 0, 1, 2$.
  \item \textbf{Reiniciar} la compleción con la regla esquemática. Esta regla no
        entra en las listas de reglas sino en el contexto, donde un
        \emph{simproc} busca la instancia $n = k$ que encaja con cada redex, la
        despliega y reescribe con ella. Así AProVE y los órdenes de términos sólo ven reglas
        de primer orden sin \isa{\textasciicircum\textasciicircum}. Los pares
        críticos de la regla esquemática se aproximan con los de sus primeras
        instancias.
\end{enumerate}

\needspace{16\baselineskip}
\begin{ejemplo}
\begin{lstlisting}
typedecl T
consts f :: "T ⇒ T"
       g :: "T ⇒ T"

complete "f (g (f x)) = f (g x)"
  sorry
(* Divergence critic: f ((g ^^ Suc n) (f x)) = f ((g ^^ Suc n) x)
   Restarting completion with schematic rule: ...
   Done: f ((g ^^ Suc ?n) (f ?x)) ≡ f ((g ^^ Suc ?n) ?x)
   Additions: f ((g ^^ Suc ?n) (f ?x)) = f ((g ^^ Suc ?n) ?x) *)
\end{lstlisting}
Sin el crítico, esta compleción no termina. Con él termina en unos segundos y deja
en \isa{compl} la regla esquemática. Las reglas esquemáticas de \isa{compl}, y la
de \cmd{complete} si ya es esquemática, pasan al contexto de las compleciones
siguientes.
\end{ejemplo}

\paragraph{Método de prueba del crítico.}
Por omisión, el crítico demuestra cada conjetura con una táctica predefinida que
hace inducción sobre el exponente, con las demás variables arbitrarias, y después
aplica simplificación y \isa{metis}, ambas con las reglas de la compleción en
curso.
Hay dos maneras de cambiarlo:
\begin{itemize}[nosep]
  \item \textbf{Añadir lemas} a la táctica predefinida con el atributo
        \isa{critic\_simps}, por ejemplo \isa{declare funpow\_swap1 [critic\_simps]}.
        Los lemas se usan en la simplificación y en \isa{metis}.
  \item \textbf{Sustituir el método} con el comando
        \begin{lstlisting}
divergence_critic_method ‹método›
\end{lstlisting}
        donde \emph{método} es cualquier método de Isar, con sus combinadores
        (\isa{,} \isa{;} \isa{|} \isa{?} \isa{+}). El texto se comprueba al
        declararlo (un método o un lema inexistente da error en ese momento) y se
        guarda en el contexto de la teoría, no en una variable global. Un texto
        vacío, \isa{divergence\_critic\_method ‹›}, vuelve a la táctica
        predefinida.
\end{itemize}
El método recibe como meta la conjetura, con sus variables libres fijadas. Como
sus nombres (\isa{n}, \isa{x}, \isa{xa}, \ldots) no se conocen de antemano,
hay dos ayudas:
\begin{itemize}[nosep]
  \item el método \cmd{exponent\_induct} hace inducción sobre la variable natural
        que aparece en los exponentes de \isa{\textasciicircum\textasciicircum}
        (debe ser única), con todas las demás arbitrarias; también sirve en
        pruebas escritas a mano;
  \item el hecho \isa{critic\_rules} contiene las reglas de la compleción en el
        momento de la conjetura, como igualdades de HOL, y sólo existe mientras el
        crítico intenta la prueba.
\end{itemize}
Cada intento tiene un límite de \isa{max\_time\_in\_critic\_proof} segundos. Si
falla, la traza lo indica con el método usado
(\isa{Schematic rule not proved (\ldots)}). Cualquier método es seguro, porque
el resultado pasa por el núcleo de Isabelle y el método sólo decide qué
conjeturas se aceptan.

\needspace{20\baselineskip}
\begin{ejemplo}
Con el método predefinido, la compleción del monoide de trenzas
$\langle a, b \mid aba = bab\rangle$ no termina. El crítico encuentra la familia,
pero no la demuestra porque el paso de inducción necesita mover una copia de
\isa{b} de la cima de la torre a su base. Cualquiera de estas dos líneas lo
resuelve:
\begin{lstlisting}
declare funpow_swap1 [critic_simps]
(* o bien *)
divergence_critic_method
  ‹exponent_induct; (simp_all add: funpow_swap1)?; metis critic_rules›

complete "a (b (a x)) = b (a (b x))" by (rule ab)
(* Done: a (b (a ?x)) ≡ b (a (b ?x)),
         a ((b ^^ Suc (Suc ?n)) (a (b ?x)))
           ≡ b (a (b (b ((a ^^ Suc ?n) ?x)))) *)
\end{lstlisting}
La compleción termina en unos 2\,s, y el resultado es exactamente el sistema
convergente infinito de la literatura,
$\{aba \rw bab\} \cup \{ab^mab \rw babba^{m-1} \mid m \ge 2\}$, escrito como
dos reglas.
\end{ejemplo}

Es un prototipo, con estas limitaciones:
\begin{itemize}[nosep]
  \item Sólo reconoce torres de un símbolo unario (\isa{Const} o \isa{Free} de tipo
        $\tau \Rightarrow \tau$) que crecen de uno en uno, con una sola variable de
        exponente.
  \item Los pares críticos de la regla esquemática con las demás se aproximan con
        sus instancias $n = 0, 1, 2$. Unificar con exponentes variables
        (Comon~\cite{comon95}) sería lo completo.
  \item La terminación de la familia se comprueba con esas tres instancias, no con
        la familia entera.
  \item Ningún método general basta, ya que Sattler-Klein~\cite{sattler-klein}
        mostró que los sistemas infinitos que produce una compleción divergente pueden
        ser muy complejos.
  \item El simplificador normal de Isabelle no aplica la regla guardada en
        \isa{compl} a términos sin \isa{\textasciicircum\textasciicircum}; el simproc
        sólo existe dentro de la compleción.
\end{itemize}

% ---------------------------------------------------------------------------
\subsection{\texttt{rec\_complete} (desactivado)}

En \isa{commands.ML} queda comentada una versión anterior,
\isa{rec\_complete "ecuación"}, que no pedía prueba. Llamaba a una ``meta-inducción''
(\isa{Proof\_Tools.meta\_induction}) que buscaba lemas y completaba en un solo paso.
La palabra clave también está comentada en \isa{IsaLibs.thy}, así que el comando no
existe en la versión actual. Lo equivalente hoy es
\isa{complete "ecuación" by inductive\_prover}.

% ===========================================================================
\section{Métodos de prueba}

% ---------------------------------------------------------------------------
\subsection{\texttt{induct\_auto}}
\label{sec:induct_auto}

\paragraph{Uso.} \isa{by induct\_auto} o \isa{apply induct\_auto}. No tiene argumentos.

\paragraph{Qué hace} (\isa{induct\_tacs4.ML}).
Para cada submeta:
\begin{enumerate}[nosep]
  \item Añade al simplificador las reglas de \isa{compl} que mencionan constantes de
        la clausura definicional de la meta (\isa{global\_simps}).
  \item Hace una búsqueda \textbf{A*} (\isa{ASTAR}). El objetivo es un estado sin
        submetas, y el coste de un estado es el número de submetas pendientes más la
        profundidad, para preferir pruebas cortas.
  \item Cada paso de la búsqueda es (inducción seguida de \isa{auto}) o sólo \isa{auto}, y debe cambiar el estado
        (\isa{CHANGED}). Las inducciones posibles (\isa{choose\_induct\_tac}) son:
    \begin{itemize}[nosep]
      \item \isa{induct\_tac} sobre cada variable inductiva, libre o meta-cuantificada,
            con cada regla de inducción de su tipo;
      \item \isa{induct} con la regla \isa{f.induct} de cada aplicación $f\,\vec{x}$
            cuyos argumentos son variables libres (inducción computacional);
      \item \isa{induct} sobre una variable inductiva con cada subconjunto de las
            demás variables como \isa{arbitrary}.
    \end{itemize}
  \item Las metas ya intentadas se guardan en una red de discriminación para no
        repetir inducciones sobre metas equivalentes.
\end{enumerate}

\paragraph{Teoría.}
Es la estrategia ``inducción y luego simplificación'' de Boyer--Moore, convertida en
búsqueda sobre el espacio de variables, reglas y generalizaciones. A* con esa
heurística no es admisible en sentido estricto (una submeta puede requerir varios
pasos), pero en la práctica encuentra pruebas cortas primero.

\begin{ejemplo}
Ambas metas se prueban directamente:
\begin{lstlisting}
lemma "add x y = add y x" by induct_auto
lemma "itrev xs ys = rev xs @ ys" by induct_auto
\end{lstlisting}
\end{ejemplo}

% ---------------------------------------------------------------------------
\subsection{\texttt{inductive\_prover}}
\label{sec:inductive_prover}

\paragraph{Uso.} \isa{by inductive\_prover}. Actúa sobre la primera submeta.

Es el componente más elaborado de IsaLibs (\isa{proof\_tools.ML}). Combina un
demostrador inductivo ``atómico'' con un bucle que analiza los fallos, conjetura
lemas, los demuestra recursivamente y los integra mediante compleción. La
figura~\ref{fig:ip} resume el flujo.

\begin{figure}[t]
\centering
\begin{tikzpicture}[
  node distance=6mm and 10mm,
  box/.style={draw,rounded corners,align=center,font=\small,inner sep=4pt,text width=4.3cm},
  dec/.style={draw,diamond,aspect=2.4,align=center,font=\small,inner sep=1pt},
  arr/.style={-{Stealth[length=2.2mm]},thick}]
  \node[box] (pre) {Preprocesar la meta $P$;\\ lemas $:=$ \isa{.simps} de la clausura};
  \node[box,below=of pre] (round) {Ronda en paralelo:\\ (a) probar $P$\\ (b) probar las conjeturas pendientes};
  \node[dec,below=of round] (res) {¿resultado?};
  \node[box,right=of res] (lemma) {Compleción condicional\\ de lemas $\cup\{\ell\}$ con\\ el lema nuevo $\ell$};
  \node[box,left=of res] (done) {Repetir la prueba\\ de $P$ con los\\ lemas hallados};
  \node[box,below=of res] (fail) {Submetas atascadas};
  \node[box,below=of fail] (gen) {Generar conjeturas (fallos,\\ exploración, Walsh, transversales)};
  \node[box,below=of gen] (filt) {Generalizar, normalizar,\\ orientar, refutar, subsunción};
  \draw[arr] (pre) -- (round);
  \draw[arr] (round) -- (res);
  \draw[arr] (res) -- node[above,font=\scriptsize]{lema} (lemma);
  \draw[arr] (res) -- node[above,font=\scriptsize]{$P$} (done);
  \draw[arr] (res) -- node[right,font=\scriptsize]{fallo} (fail);
  \draw[arr] (fail) -- (gen);
  \draw[arr] (gen) -- (filt);
  \draw[arr] (lemma.north) |- (round.east);
  \draw[arr] (filt.west) -- ($(done.west |- filt.west) + (-4mm,0)$) |- (round.west);
\end{tikzpicture}
\caption{Bucle principal de \cmd{inductive\_prover} (\isa{recursive\_prover}).}
\label{fig:ip}
\end{figure}

\paragraph{1. Preparación.}
\begin{itemize}[nosep]
  \item La meta se \emph{rulifica} (cuantificadores e implicaciones de HOL pasan a
        $\bigwedge$ y $\Longrightarrow$).
  \item Los \emph{lemas iniciales} son las ecuaciones \isa{.simps} de la clausura
        definicional de la meta (\isa{get\_definitional\_rewrites}). Se quitan del
        simplificador estándar (\isa{ctxt\_nodefs}) para que sólo se usen bajo el
        control de la terminación.
  \item La traza muestra la meta en \LaTeX{} (\isa{Latex Goal}) y en texto.
  \item Se fija el oráculo \isa{terminates = Aprove.memoized\_terminates}, y
        \isa{safe\_rules} selecciona la parte terminante de un conjunto de lemas
        (más las ecuaciones permutativas) con
        \isa{Conditional\_Completion.discriminate\_object\_rules}.
\end{itemize}

\paragraph{2. Demostrador atómico} (\isa{Inductive\_Tacs.induct\_auto\_failure\_tac}).
Es una búsqueda A* con el mismo coste que \cmd{induct\_auto}. Cada paso aplica a
todas las submetas \isa{TRY} (inducción) seguido de \isa{TRY auto}, con el
simplificador extendido con \isa{safe\_rules(lemas)}. La inducción puede ser:
\begin{itemize}[nosep]
  \item \emph{computacional}, sobre cada aplicación $f\,\vec x$ con argumentos libres,
        de menor a mayor, con y sin \isa{arbitrary} para el resto de variables;
  \item \emph{estructural}, sobre cada variable libre de tipo inductivo, también con
        y sin \isa{arbitrary}.
\end{itemize}
Cuando una submeta pendiente ya no tiene variables libres (todas quedaron
meta-cuantificadas tras la inducción), se \textbf{registra como fallo} en una red
(\isa{check\_fail\_proof\_state}). Esas submetas son las candidatas naturales a lema.
Cada intento tiene un límite de \isa{max\_time\_in\_proof} segundos.

\paragraph{3. Ronda} (\isa{next\_subgoal}).
Con el conjunto actual de lemas se lanzan dos futuros en paralelo:
\begin{enumerate}[label=(\alph*),nosep]
  \item normalizar la meta original con las reglas seguras; si queda en
        \isa{True}, o si el demostrador atómico la prueba, se lanza la excepción
        \isa{DONE}, que cancela todo lo pendiente;
  \item normalizar, orientar, deduplicar y filtrar las conjeturas pendientes, y
        probarlas todas en paralelo (\isa{Par\_List.get\_some}), de modo que gana
        la primera que se demuestra.
\end{enumerate}
Si en la ronda anterior no se añadió ningún lema, (a) no se repite.

\paragraph{4. Integración de un lema.}
Si el lema contiene $\bigwedge$ o $\Longrightarrow$ (\isa{unsupported\_rule}), pasa a
una lista aparte (\isa{mlemmas}), y el demostrador atómico lo usa como regla, pero no
entra en la compleción. En otro caso se recalcula el conjunto de lemas con
\textbf{compleción condicional} (\isa{build\_trs}, que llama a
\isa{Conditional\_Completion.run\_completion}). Esta variante de la
sección~\ref{sec:complete} admite reglas con premisas:
\begin{itemize}[nosep]
  \item las premisas que no son ecuaciones se convierten en
        $P = \mathit{True}$ o $P = \mathit{False}$;
  \item cada par crítico entre reglas condicionales se demuestra explícitamente con
        \isa{frule} + \isa{frule} + \isa{auto};
  \item la precedencia sigue el orden de dependencias definicionales;
  \item si ninguna orientación del lema nuevo es terminante y el lema no es
        permutativo, el conjunto de lemas no cambia.
\end{itemize}
El lema y los ya probados se guardan en la red de subsunción, y se pasa a una
nueva ronda.

\paragraph{5. Generación de conjeturas tras un fallo.}
Las conjeturas que no se pudieron probar se añaden a la red de ``ya intentadas''. Con
las submetas atascadas se construyen cuatro secuencias perezosas, que se intercalan
de forma equitativa (\isa{Seq.interleave}):
\begin{description}[nosep,font=\normalfont\ttfamily]
  \item[EQ\_FROM\_FAILURE] las submetas atascadas mismas, ordenadas de más general
        a menos general;
  \item[TX\_ON\_FAILURES] exploración de teorías sobre los subtérminos (y sus
        generalizaciones) de cada fallo
        (\isa{equational\_theory\_exploration\_on\_conj});
  \item[WALSH] el crítico de divergencia aplicado a pares (submeta anterior, fallo).
        Busca con unificación de patrones $G$ y $H$ tales que $s = G\,U$ y
        $s' = G\,(H\,U)$, y propone $G(H\,U_0) = H(G\,U_0)$ y
        $G(H\,U_0) = G(G\,U_0)$. La variante condicional sólo actúa si ambas
        submetas tienen las mismas premisas, y las conserva;
  \item[TRANSVERSE] reglas de onda transversales para acumuladores. Unifica
        $s = G\,U\,\mathit{Acc}$ y $s' = G\,(H\,U)\,\mathit{Acc}$, propone la meta
        esquemática $G\,(H\,U_0)\,\mathit{Acc} = G\,U_0\,(?F\,\mathit{Acc})$ y la
        intenta demostrar con el demostrador atómico
        (\isa{Prover.prove\_schematic\_from\_method}). Si lo logra, $?F$ queda
        instanciada y la conjetura resultante es concreta.
\end{description}
Antes de empezar, el almacén de conjeturas ya contiene otras dos secuencias, que se
consumen antes que las generadas por los fallos porque las nuevas se añaden al final
(\isa{add\_sequence}):
\begin{description}[nosep,font=\normalfont\ttfamily]
  \item[EQ\_KNOWN] (si \isa{schematic\_conjectures}) los esquemas de
        \isa{prop\_scheme} instanciados con las constantes del problema y no
        refutados (sección~\ref{sec:prop_scheme});
  \item[EQ\_TX] (si \isa{theory\_exploration\_conjectures}) exploración de teorías
        sobre toda la signatura de la meta (ver abajo).
\end{description}

\paragraph{6. Filtro y generalización de cada candidata.}
Cada conjetura pasa en paralelo por \isa{generalize1} y \isa{generalize2}. Ambas hacen
lo mismo en distinto orden:
\begin{itemize}[nosep]
  \item Para \textbf{generalizar hipótesis} (\isa{generalize\_assumptions}) se elige
        el subconjunto más pequeño de premisas, ordenadas por tamaño, con el que la
        conjetura no se refuta.
  \item Para \textbf{generalizar subtérminos comunes} (\isa{generalize\_cond\_eq}),
        los subtérminos de tipo base que aparecen más de una vez se sustituyen por
        variables nuevas. La generalización se conserva sólo si no se refuta.
  \item Para \textbf{simplificar} (\isa{gen\_simplify}) se siguen estos pasos:
    \begin{itemize}[nosep]
      \item normalizar con las reglas seguras;
      \item aplicar las heurísticas de \isa{apply\_heuristics}, que cambian
            $C\,x = C\,y$ por $x = y$ si $C$ es un constructor, y
            $(\lambda x.\,s) = (\lambda y.\,t)$ por $s = t$ (extensionalidad);
      \item orientar;
      \item descartar si es trivial o si está subsumida por algo ya intentado.
    \end{itemize}
\end{itemize}
En cada ronda se toman tantas conjeturas como procesadores
(\isa{Utils.num\_processors}). Se rechazan las que igualan funciones y las que sólo
contienen conectivos lógicos y constructores (\isa{filter\_conjectures}).

\paragraph{7. Exploración de teorías} (\isa{equivalence\_terms.ML}).
Es una versión de QuickSpec~\cite{quickspec} construida sobre Isabelle:
\begin{enumerate}[nosep]
  \item Para los \textbf{datos de prueba}, los tipos de la clausura definicional se
        \emph{monomorfizan}. Las variables de tipo se sustituyen por un tipo no
        paramétrico del problema o, si no lo hay, por el tipo auxiliar
        \isa{datatype two = A | B} de \isa{IsaLibs.thy}. Para cada tipo se enumeran
        \isa{max\_random\_terms} valores cerrados
        (\isa{Counter\_Example.enumerate\_terms}).
  \item Para obtener los \textbf{términos}, con las constantes de la clausura y
        \isa{max\_vars\_in\_tx} variables por tipo se enumeran \emph{todos} los $\lambda$-términos en forma
        $\beta$-normal del tipo $\tau_1 \Rightarrow \cdots \Rightarrow \tau_n
        \Rightarrow \sigma$, hasta tamaño \isa{max\_lambda\_size}$+n$
        (\isa{Random\_Terms.enumerate\_all\_lambda\_terms}, basado en el conteo de
        términos de Zaionc). Se aplican a la signatura, de modo que cada uno
        representa un término sobre las funciones del problema, y se ordenan de más
        general a menos general.
  \item Para formar las \textbf{clases de equivalencia}, cada término se normaliza
        con las reglas seguras. Su \emph{huella} es la lista de formas normales de sus instancias
        con los valores de prueba. Dos términos con la misma huella generan la
        conjetura $t = s$ (bajo las premisas de la meta), que se normaliza y se
        filtra con Quickcheck y Nitpick.
\end{enumerate}

\paragraph{8. Final.}
Cuando la meta original se demuestra, la traza muestra
``\isa{Proved original goal!}'' y el conjunto de lemas hallados (\isa{Simpset}, también
en \LaTeX). El método vuelve a probar la meta con el demostrador atómico usando esos
lemas, y es esa prueba la que produce el teorema. Todo el proceso tiene el límite
\isa{max\_time\_in\_inductive\_proof}.

\begin{ejemplo}
Tiempos medidos en Isabelle2026 con AProVE (8 hilos):
\begin{center}
\begin{tabular}{@{}lrl@{}}
\toprule
Meta & Tiempo & Lemas descubiertos (además de \isa{.simps})\\
\midrule
\isa{add x y = add y x} & 37\,s & \isa{add x Z = x}, \isa{add x (S y) = S (add x y)}\\
\isa{add (add x y) z = add x (add y z)} & 26\,s & ninguno\\
\isa{rev2 (rev2 xs) = xs} & 316\,s & \isa{app x [] = x}, asociatividad de \isa{app},\\
 & & \isa{rev2 (app xs ys) = app (rev2 ys) (rev2 xs)}\\
\bottomrule
\end{tabular}
\end{center}
\end{ejemplo}

\paragraph{Opciones relevantes.}
\isa{max\_time\_in\_inductive\_proof}, \isa{max\_time\_in\_proof},
\isa{max\_time\_in\_counter\_ex}, \isa{schematic\_conjectures},
\isa{theory\_exploration\_conjectures}, \isa{max\_lambda\_size},
\isa{max\_random\_terms}, \isa{max\_vars\_in\_tx}, \isa{max\_time\_normalization} y
las de terminación.

% ---------------------------------------------------------------------------
\subsection{\texttt{inductive\_sledgehammer\_prover}}

Es el mismo bucle de la sección~\ref{sec:inductive_prover}, con otro demostrador
atómico (\isa{induct\_auto\_sledgehammer\_failure\_tac}):
\begin{itemize}[nosep]
  \item los lemas se registran como hechos del contexto (\isa{note\_thmss}) en lugar
        de reglas de simplificación;
  \item tras cada paso de inducción y \isa{auto}, a las submetas cerradas se les
        aplica Sledgehammer con reconstrucción \isa{metis}
        (\isa{Sledgehammer\_Tactics.sledgehammer\_with\_metis\_tac}), con el límite
        de la opción del sistema \isa{sledgehammer\_timeout}, que \isa{IsaLibs.thy}
        fija en 6\,s;
  \item las submetas que siguen pendientes se registran como fallos.
\end{itemize}
Sirve cuando el lema necesario se deduce de hechos de la biblioteca de HOL que el
simplificador no encuentra. En Isabelle2026, \isa{add x y = add y x} se demuestra en
212\,s con este método. El conjunto final de lemas contiene \isa{add x Z = x} y
\isa{add x (S y) = S (add x y)}. Durante la búsqueda también demuestra la
asociatividad de \isa{add}, aunque ésta no queda en ese conjunto.

% ---------------------------------------------------------------------------
\subsection{\texttt{inductive\_provers}}

Es el mismo bucle con el demostrador atómico \isa{induct\_failure\_provers\_tac},
que intenta primero \isa{induct\_auto\_failure\_tac} y, sólo si no resuelve la meta
(\isa{SOLVE \ldots\ ORELSE SOLVE \ldots}), la variante con Sledgehammer. Combina la
rapidez de la primera con la fuerza de la segunda. Por ejemplo,
\isa{add (add x y) z = add x (add y z)} se demuestra en 27\,s.

% ---------------------------------------------------------------------------
\subsection{\texttt{spass\_auto}}

Llama a \isa{Sledgehammer\_Tactics.sledgehammer\_with\_metis\_tac} sobre la primera
submeta, con un límite de \isa{max\_time\_in\_proof} segundos y los demostradores
por omisión de Sledgehammer, y reconstruye la prueba con \isa{metis}. El nombre es
histórico, porque cuando se escribió SPASS era el demostrador por omisión. Internamente
lo usa \isa{Prover.prove\_from\_spass}.
\begin{lstlisting}
lemma "rev (xs @ ys) = rev ys @ rev xs"
  by spass_auto
\end{lstlisting}

\begin{nota}
En Isabelle2026, Isar evalúa cada método primero sobre un estado ficticio sin
submetas (\isa{TERM \_}). Ahí \isa{sledgehammer\_with\_metis\_tac} lanza el error
``\isa{Subgoal number 1 out of range}'' en lugar de fallar en silencio, y el error
abortaba todo el comando. Por eso la llamada va ahora dentro de \isa{SUBGOAL}. Los
demostradores atómicos con Sledgehammer ya la llamaban dentro de \isa{SUBGOAL}, así
que no tenían el problema.
\end{nota}

% ---------------------------------------------------------------------------
\subsection{\texttt{induct\_auto\_sledgehammer\_failure\_tac}}

Expone como método el demostrador atómico de
\cmd{inductive\_sledgehammer\_prover}, sin conjunto de lemas y descartando los
fallos registrados. Es útil para depurar, ya que muestra qué puede probar el paso
atómico por sí solo.

% ===========================================================================
\section{Atributos y conjuntos de teoremas}

\subsection{\texttt{compl}}

Es un conjunto con nombre (\isa{Named\_Thms}, estructura \isa{Completion\_Rules}),
con atributo \isa{[compl]} y \isa{[compl del]}. Contiene el sistema de reescritura
que va construyendo \cmd{complete}. Lo usan:
\begin{itemize}[nosep]
  \item \cmd{complete}, como parte de $E$ (sólo las reglas que mencionan constantes
        de la nueva ecuación);
  \item \cmd{induct\_auto}, que añade al simplificador las reglas relevantes.
\end{itemize}
Se consulta con \isa{thm compl}. Se le pueden añadir teoremas a mano con
\isa{declare foo [compl]}, aunque así no pasan por la compleción.

\subsection{\texttt{prop\_scheme}}
\label{sec:prop_scheme}

Es un conjunto con nombre (\isa{PropSchemes}) de \emph{esquemas de propiedades}.
Cada esquema es una definición cuyos parámetros son funciones:
\begin{lstlisting}
definition associativity where
  [prop_scheme]: "associativity R ≡ ∀x y z. R (R x y) z = R x (R y z)"
\end{lstlisting}
\isa{IsaLibs.thy} predefine los esquemas
\isa{associativity}, \isa{commutativity}, \isa{associativity\_commutativity},
\isa{left\_identity}, \isa{right\_identity}, \isa{absorbing\_element},
\isa{left\_absorbing\_element}, \isa{right\_absorbing\_element}, \isa{idempotent},
\isa{distributive}, \isa{left\_distributive}, \isa{right\_distributive},
\isa{arguments\_output}, \isa{injectivity\_one}, \isa{injectivity\_two} e
\isa{injectivity\_three}.

\paragraph{Uso.} Con \isa{schematic\_conjectures = true}, que es el valor en
\isa{IsaLibs.thy}, \cmd{inductive\_prover} llama a \isa{DB\_EQ\_Terms.known\_laws}:
\begin{enumerate}[nosep]
  \item instancia cada parámetro de cada esquema con todas las constantes de la
        clausura definicional de tipo compatible
        (\isa{Inst\_Utils.compatible\_arguments}, \isa{Inst\_Utils.combine});
  \item desdobla la definición del esquema y elimina el $\forall$ exterior;
  \item descarta \isa{True}, \isa{False} y lo que Quickcheck o Nitpick refutan.
\end{enumerate}
\begin{ejemplo}
Para la meta \isa{add x y = add y x}, las conjeturas que sobreviven son:
\begin{lstlisting}
add x Z = x                           (right_identity)
add x y = add y x                     (commutativity)
(add x = add y) = (x = y)             (injectivity_one)
add (add x y) z = add x (add y z)     (associativity)
add x (add y z) = add y (add x z)     (associativity_commutativity)
\end{lstlisting}
\end{ejemplo}
La idea central de IsaScheme~\cite{montano2012} es que el conocimiento matemático
general entra como plantillas de orden superior. Un usuario puede añadir sus propios
esquemas con el atributo.

% ===========================================================================
\section{Opciones de configuración}

Todas son opciones de contexto de Isabelle (\isa{Attrib.config\_*}) y se cambian con
\isa{declare [[opción = valor]]}. La tabla da el valor por omisión del código ML y
el valor que fija \isa{IsaLibs.thy}, que es el que rige en las teorías de usuario.

\begin{longtable}{@{}p{4.9cm}ccp{6.4cm}@{}}
\toprule
\textbf{Opción} & \textbf{ML} & \textbf{IsaLibs} & \textbf{Significado}\\
\midrule
\endhead
\isa{use\_aprove} & true & true & Usar AProVE. Con \isa{false}, toda prueba de terminación falla.\\
\isa{use\_aprove\_srv} & false & \textbf{true} & Hablar con el servidor de AProVE vía \isa{nc} en lugar de ejecutar el \isa{jar}.\\
\isa{max\_time\_in\_termination} & 180 & 20 & Segundos para AProVE (\isa{-t}) y para el método de pares de dependencia interno.\\
\isa{generate\_cps} & true & true & Deducir pares críticos en la compleción (también en la condicional de \cmd{inductive\_prover}).\\
\isa{divergence\_critic} & true & true & Con \isa{generate\_cps}, generalizar familias divergentes de reglas con \isa{\textasciicircum\textasciicircum} (prototipo).\\
\isa{max\_time\_in\_critic\_proof} & 60 & 60 & Segundos para cada prueba del crítico de divergencia.\\
\isa{use\_quickcheck} & true & true & Buscar contraejemplos con Quickcheck.\\
\isa{use\_nitpick} & true & true & Buscar contraejemplos con Nitpick.\\
\isa{simp\_before} & false & false & Simplificar la conjetura antes de buscar contraejemplos.\\
\isa{max\_time\_in\_counter\_ex} & 5 & 25 & Segundos por búsqueda de contraejemplo.\\
\isa{max\_time\_in\_proof} & 10 & 60 & Segundos por intento del demostrador atómico y para \cmd{spass\_auto}.\\
\isa{max\_time\_in\_inductive\_proof} & 3600 & 7200 & Límite global de \cmd{inductive\_prover}.\\
\isa{max\_time\_normalization} & 5 & 5 & Límite de cada normalización con el simplificador.\\
\isa{schematic\_conjectures} & false & \textbf{true} & Conjeturas por esquemas \isa{prop\_scheme} (\isa{EQ\_KNOWN}).\\
\isa{theory\_exploration\_conjectures} & false & \textbf{true} & Conjeturas por exploración de teorías (\isa{EQ\_TX}).\\
\isa{max\_lambda\_size} & 10 & 10 & Tamaño máximo de los $\lambda$-términos enumerados.\\
\isa{max\_random\_terms} & 10 & 8 & Valores de prueba por tipo.\\
\isa{max\_vars\_in\_tx} & 3 & 3 & Variables por tipo en la exploración.\\
\isa{max\_num\_generalizations} & 2 & 2 & Sin efecto, porque sólo lo usa código comentado.\\
\isa{max\_time\_in\_fitness} & 30 & 15 & Programación genética (sólo ML).\\
\bottomrule
\end{longtable}

\isa{IsaLibs.thy} también cambia opciones estándar de Isabelle:
\begin{itemize}[nosep]
  \item \isa{quick\_and\_dirty = true}, así que \isa{sorry} funciona en todas las
        teorías que importan IsaLibs;
  \item \isa{quickcheck\_quiet = true};
  \item \isa{linarith\_split\_limit = 10};
  \item \isa{eta\_contract = false};
  \item la opción del sistema \isa{sledgehammer\_timeout = 6}, sólo dentro del proceso.
\end{itemize}

% ===========================================================================
\section{Antiquotaciones ML}

\isa{utils.ML} define dos antiquotaciones para escribir patrones en código ML:
\begin{lstlisting}
val T = @{typ_pat "?'a ⇒ ?'b"}
val P = @{term_pat "?X (?lhs :: ?'a) ≡ ?X (?rhs :: ?'a)"}
\end{lstlisting}
Leen el texto en modo esquemático, por lo que admiten \isa{?x} y \isa{?'a}, que
\isa{@\{term\}} rechazaría. Se usan en los críticos de divergencia y en
\isa{heuristic1}.

% ===========================================================================
\section{Mapa de módulos}

\begin{longtable}{@{}p{4.3cm}p{10.4cm}@{}}
\toprule
\textbf{Archivo} & \textbf{Contenido}\\
\midrule
\endhead
\isa{IsaLibs.thy} & Palabras clave, carga de módulos, opciones, \cmd{orient\_rules} y esquemas predefinidos.\\
\isa{commands.ML} & Comando \cmd{complete} (y \isa{rec\_complete} comentado).\\
\isa{oriented\_rules.ML} & Comando \cmd{orient\_rules}.\\
\isa{schematic\_lemmas.ML} & Comando \cmd{schematic\_lemma}.\\
\isa{tables.ML} & Tablas auxiliares y el conjunto \isa{compl}.\\
\isa{utils.ML} & Utilidades generales de orientación, normalización, dependencias definicionales, plazos (\isa{apply\_timeout}), HTML y \LaTeX, y antiquotaciones.\\
\isa{ground\_completion.ML} & Compleción sin fallo con ecuaciones incondicionales, reglas esquemáticas y crítico de divergencia (con el comando \cmd{divergence\_critic\_method} y el método \cmd{exponent\_induct}).\\
\isa{conditional\_completion.ML} & Compleción con reglas condicionales (la usa \cmd{inductive\_prover}).\\
\isa{orders.ML} & Extensiones de precedencias y LPO para unibilidad en términos cerrados.\\
\isa{aprove.ML} & Traducción a WST e interfaz con AProVE.\\
\isa{dependency\_pair.ML}, \isa{sat\_interface.ML} & Pares de dependencia estáticos y codificación SAT (no se usan en la ruta principal).\\
\isa{prover.ML} & Ejecución de métodos con plazo, pruebas esquemáticas, \cmd{spass\_auto}.\\
\isa{induct\_tacs4.ML} & Método \cmd{induct\_auto}.\\
\isa{induct\_tacs.ML} & Demostradores atómicos con registro de fallos.\\
\isa{proof\_tools.ML} & \cmd{inductive\_prover} y sus variantes.\\
\isa{divergence.ML} & Crítico de divergencia de Walsh y reglas de onda transversales.\\
\isa{equivalence\_terms.ML} & Exploración de teorías y conjeturas por esquemas.\\
\isa{random\_terms.ML}, \isa{enumerated\_terms.ML}, \isa{counting\_terms.ML} & Conteo, enumeración y generación aleatoria de $\lambda$-términos tipados.\\
\isa{counter\_example.ML} & Quickcheck, Nitpick y su memorización.\\
\isa{inst\_utils.ML} & Instanciación de esquemas.\\
\isa{meta\_gp\_hol.ML}, \isa{exhaust.ML} & Programación genética sobre $\lambda$-cálculo (estructura \isa{GP}).\\
\isa{smt.ML} & Lectura y conversión de problemas TIP en SMT-LIB.\\
\isa{latex.ML}, \isa{gnuplot.ML}, \isa{mysql.ML} & Salida de resultados.\\
\isa{papers/} & Arneses de experimentos (\isa{ESWA2016.ML}, cruce rectilíneo, base de datos).\\
\isa{gp.ML}, \isa{gp\_hol.ML}, \isa{html.ML}, \isa{induct\_search.ML}, \isa{def\_utils.ML} & No se cargan (código antiguo).\\
\bottomrule
\end{longtable}

% ===========================================================================
\section{Advertencias prácticas}

\begin{enumerate}
  \item \textbf{AProVE es imprescindible.} Sin él (o sin su servidor, con
        \isa{use\_aprove\_srv = true}) ninguna regla pasa la prueba de
        terminación. El síntoma es que \cmd{inductive\_prover} falla siempre y
        \cmd{complete} no orienta nada nuevo. Lo mismo pasa si el servidor queda
        bloqueado; la sección~\ref{sec:carga} explica cómo comprobarlo y reiniciarlo.
  \item \textbf{Estado global.} Las tablas de \cmd{orient\_rules},
        \cmd{schematic\_lemma}, terminación, contraejemplos, discriminación de reglas
        y unibilidad son variables ML globales. Duran todo el proceso, quedan en el
        heap y no siguen el contexto de la teoría. AProVE recibe los nombres
        \emph{base} de las constantes, pero la clave de la memoria es el término
        completo, así que no se confunden constantes homónimas de teorías distintas.
        Por la misma razón, estas tablas guardan términos o posiciones, nunca
        teoremas. jEdit vuelve a ejecutar los comandos sobre versiones nuevas de la
        teoría, y un teorema guardado de una ejecución anterior no se puede
        transferir a la nueva. Antes, la tabla de \isa{discriminate\_rules} de la
        compleción básica guardaba teoremas, y \cmd{complete} fallaba con
        ``\isa{Cannot transfer: not a super theory}''. Ahora guarda posiciones,
        igual que la de la compleción condicional.
  \item \textbf{\isa{quick\_and\_dirty}} está activado. Antes de dar por buena una
        teoría conviene comprobar que no quedan \isa{sorry}.
  \item \textbf{Pares críticos activados} por omisión
        (\isa{generate\_cps = true}). Con ellos la compleción puede no terminar.
        El crítico de divergencia resuelve algunos de esos casos, y
        \isa{declare [[generate\_cps = false]]} vuelve a la compleción sin pares
        críticos.
  \item \textbf{Paralelismo.} \cmd{inductive\_prover} usa muchos hilos (futuros,
        \isa{Par\_List}). Los tiempos dependen mucho del número de núcleos y de
        \isa{threads}.
  \item \textbf{Trazas.} El progreso se ve en el panel de \emph{tracing} de jEdit,
        con mensajes como \isa{Proving:}, \isa{Proved!:}, \isa{Normalising (EQ\_TX):},
        \isa{subsumed or rejected}, \isa{Completing:}, \isa{Done:},
        \isa{Elapsed time}. Las referencias ML \isa{Ground\_Completion.completion\_debug}
        y similares activan más detalle. En \isa{isabelle build} la traza no queda en
        el registro; para capturarla se puede redirigir
        \isa{Private\_Output.tracing\_fn}.
\end{enumerate}

% ===========================================================================
\begin{thebibliography}{99}

\bibitem{montano2012}
O.~Montaño-Rivas, R.~McCasland, L.~Dixon, A.~Bundy.
\newblock Scheme-based theorem discovery and concept invention.
\newblock \emph{Expert Systems with Applications} 39(2):1637--1646, 2012.

\bibitem{baader-nipkow}
F.~Baader, T.~Nipkow.
\newblock \emph{Term Rewriting and All That}.
\newblock Cambridge University Press, 1998.

\bibitem{knuth-bendix}
D.~E.~Knuth, P.~B.~Bendix.
\newblock Simple word problems in universal algebras.
\newblock En \emph{Computational Problems in Abstract Algebra}, Pergamon, 1970.

\bibitem{bdp}
L.~Bachmair, N.~Dershowitz, D.~A.~Plaisted.
\newblock Completion without failure.
\newblock En \emph{Resolution of Equations in Algebraic Structures}, vol.~2,
  Academic Press, 1989.

\bibitem{martin-nipkow}
U.~Martin, T.~Nipkow.
\newblock Ordered rewriting and confluence.
\newblock En \emph{CADE-10}, LNCS 449, Springer, 1990.

\bibitem{miller}
D.~Miller.
\newblock A logic programming language with lambda-abstraction, function
  variables, and simple unification.
\newblock \emph{Journal of Logic and Computation} 1(4), 1991.

\bibitem{arts-giesl}
T.~Arts, J.~Giesl.
\newblock Termination of term rewriting using dependency pairs.
\newblock \emph{Theoretical Computer Science} 236, 2000.

\bibitem{aprove}
J.~Giesl et~al.
\newblock Analyzing program termination and complexity automatically with AProVE.
\newblock \emph{Journal of Automated Reasoning} 58(1), 2017.

\bibitem{kusakari}
K.~Kusakari, Y.~Isogai, M.~Sakai, F.~Blanqui.
\newblock Static dependency pair method based on strong computability for
  higher-order rewrite systems.
\newblock \emph{IEICE Transactions on Information and Systems} E92-D(10), 2009.

\bibitem{boyer-moore}
R.~S.~Boyer, J.~S.~Moore.
\newblock \emph{A Computational Logic}.
\newblock Academic Press, 1979.

\bibitem{ireland-bundy}
A.~Ireland, A.~Bundy.
\newblock Productive use of failure in inductive proof.
\newblock \emph{Journal of Automated Reasoning} 16, 1996.

\bibitem{walsh}
T.~Walsh.
\newblock A divergence critic for inductive proof.
\newblock \emph{Journal of Artificial Intelligence Research} 4, 1996.

\bibitem{kirchner89}
H.~Kirchner.
\newblock Schematization of infinite sets of rewrite rules generated by divergent
  completion processes.
\newblock \emph{Theoretical Computer Science} 67(2--3), 1989.

\bibitem{chen-hsiang-kong}
H.~Chen, J.~Hsiang, H.-C.~Kong.
\newblock On finite representations of infinite sequences of terms.
\newblock En \emph{CTRS 1990}, LNCS 516, Springer, 1991.

\bibitem{comon95}
H.~Comon.
\newblock On unification of terms with integer exponents.
\newblock \emph{Mathematical Systems Theory} 28, 1995.

\bibitem{thomas-watson}
M.~Thomas, P.~Watson.
\newblock Solving divergence in Knuth--Bendix completion by enriching signatures.
\newblock \emph{Theoretical Computer Science} 112, 1993.

\bibitem{sattler-klein}
A.~Sattler-Klein.
\newblock Divergence phenomena during completion.
\newblock En \emph{RTA 1991}, LNCS 488, Springer, 1991.

\bibitem{rippling}
A.~Bundy, D.~Basin, D.~Hutter, A.~Ireland.
\newblock \emph{Rippling: Meta-Level Guidance for Mathematical Reasoning}.
\newblock Cambridge University Press, 2005.

\bibitem{quickspec}
K.~Claessen, N.~Smallbone, J.~Hughes.
\newblock QuickSpec: guessing formal specifications using testing.
\newblock En \emph{TAP 2010}, LNCS 6143, Springer, 2010.

\bibitem{isacosy}
M.~Johansson, L.~Dixon, A.~Bundy.
\newblock Conjecture synthesis for inductive theories.
\newblock \emph{Journal of Automated Reasoning} 47(3), 2011.

\bibitem{quickcheck}
L.~Bulwahn.
\newblock The new Quickcheck for Isabelle: random, exhaustive and symbolic testing
  under one roof.
\newblock En \emph{CPP 2012}, LNCS 7679, Springer, 2012.

\bibitem{nitpick}
J.~C.~Blanchette, T.~Nipkow.
\newblock Nitpick: a counterexample generator for higher-order logic based on a
  relational model finder.
\newblock En \emph{ITP 2010}, LNCS 6172, Springer, 2010.

\end{thebibliography}

\end{document}
